\documentclass[11pt]{article}

\usepackage[T1]{fontenc}
\usepackage[utf8]{inputenc}
\usepackage{amsmath,amssymb,amsthm}
\usepackage[margin=1in]{geometry}
\usepackage[hidelinks,pdftitle={The Center and Ozone Group of Uq+(B2) at a Primitive Cube Root},pdfauthor={Xiaojian Zeng},pdfsubject={Quantum algebra at a primitive cube root},pdfkeywords={quantum algebra, root of unity, center, Ore extension, ozone group}]{hyperref}

\newtheorem{theorem}{Theorem}[section]
\newtheorem{lemma}[theorem]{Lemma}
\newtheorem{proposition}[theorem]{Proposition}
\newtheorem{remark}[theorem]{Remark}

\newcommand{\Oz}{\operatorname{Oz}}
\newcommand{\Aut}{\operatorname{Aut}}
\newcommand{\gr}{\operatorname{gr}}
\newcommand{\id}{\mathrm{id}}
\newcommand{\kk}{\mathbf{k}}

\title{The Center and Ozone Group of $U_q^+(B_2)$\\at a Primitive Cube Root}
\author{Xiaojian Zeng\\
PBC School of Finance, Tsinghua University, 100190\\
\texttt{xiaojianzeng@foxmail.com}}
\date{}

\begin{document}

\maketitle

\begin{abstract}
Let $\kk$ be an algebraically closed field of characteristic zero and let $q\in\kk^*$
be a primitive cube root of unity.  For the positive quantum algebra
$A=U_q^+(B_2)$, we give a direct Ore-extension presentation at this
parameter, determine its full center, and prove that every automorphism fixing
the center pointwise is the identity.  More precisely, for the standard root
vectors $e_3,z$ and the central element $z_1$ defined below, we prove
\[
 Z(A)=\kk[e_1^3,e_2^3,e_3^3,z,z_1]
 \qquad\text{and}\qquad
 \Oz(A)=\{\id_A\}.
\]
The argument is specific to order three and does not obtain the conclusion by
substituting into results whose hypotheses require root order at least five.
\end{abstract}

\noindent\textbf{Keywords.}
Quantum algebra; root of unity; center; Ore extension; ozone group.

\medskip
\noindent\textbf{2020 Mathematics Subject Classification.}
16W20, 16S36, 17B37.

\section{Context and statement}

For a $\kk$-algebra $R$, its ozone group is
\[
 \Oz(R)=\Aut_{Z(R)\text{-alg}}(R),
\]
the group of algebra automorphisms fixing every element of the center.  This
terminology was introduced by Chan, Gaddis, Won, and Zhang
\cite{ChanGaddisWonZhang2024}.

The root-of-unity center of $U_q^+(B_2)$ was determined by Bera and
Mukherjee under the standing hypothesis that $q$ has primitive order
$m\geq 5$ \cite[Theorem~9.2]{BeraMukherjee2026}.  Under the same restriction,
G\'omez and Venegas computed the ozone group as
$\mu_{\gcd(\ell,2)}$, where $\ell=\operatorname{ord}(q^2)$
\cite[Theorem~3.7]{GomezVenegas2026}.  The formal value at $m=3$ is trivial,
but order three lies outside both hypotheses.  It is also a degenerate
specialization: the two middle coefficients in the second quantum Serre
relation vanish.  We therefore verify the presentation, center, and
automorphism argument directly.  The root-vector notation and a related PBW
presentation occur in the general type-$B_2$ literature
\cite[Section~3.1.2]{AndruskiewitschDumas2008}; no such presentation is used
here without verification at the cubic parameter.

Let $A$ be generated by $e_1,e_2$ with relations
\begin{align*}
 e_1^2e_2-(q^2+q^{-2})e_1e_2e_1+e_2e_1^2&=0,\\
 e_2^3e_1-(q^2+1+q^{-2})e_2^2e_1e_2
 +(q^2+1+q^{-2})e_2e_1e_2^2-e_1e_2^3&=0.
\end{align*}
Define
\begin{align}
 e_3&=e_1e_2-q^2e_2e_1,\label{eq:e3}\\
 z&=e_2e_3-q^2e_3e_2,\label{eq:z}\\
 z_1&=e_1e_2e_3+\frac{e_1z}{q^2-1}
                    +\frac{e_3^2}{q^4-1}.\label{eq:z1}
\end{align}
Both denominators in \eqref{eq:z1} are nonzero at a primitive cube root.

\begin{theorem}\label{thm:main}
Let $\kk$ be an algebraically closed field of characteristic zero, let
$q\in\kk^*$ be a primitive third root of unity, and let $A$ be exactly the
algebra above.  Then
\[
 Z(A)=\kk[e_1^3,e_2^3,e_3^3,z,z_1]
 \quad\text{and}\quad
 \Aut_{Z(A)\text{-alg}}(A)=\{\id_A\}.
\]
The center equality is an equality of generated subalgebras; no algebraic
independence of the five displayed generators is asserted.
\end{theorem}

This treats exactly the parameter excluded from the cited $m\geq5$ formulas.
It is not a claim of publication priority for the cubic case.

\section{A cubic Ore presentation}

Put $x=e_1$, $y=e_2$, and $t=e_3$.  Since
$q^3=1$, $q\neq1$, and $1+q+q^2=0$, the two defining relations specialize to
\begin{equation}
 x^2y+xyx+yx^2=0,
 \qquad
 y^3x=xy^3.
 \label{eq:specialized-serre}
\end{equation}

\begin{lemma}\label{lem:ore}
There is an isomorphism
\[
 A\cong \kk[z,t][x;\sigma_1][y;\sigma_2,\delta_2]
\]
whose commutation rules, in the PBW order $z<t<x<y$, are
\begin{equation}
 z\text{ is central},\qquad
 xt=qtx,\qquad
 yt=q^2ty+z,\qquad
 yx=qxy-qt.
 \label{eq:ore-rules}
\end{equation}
Consequently,
\[
 \{z^dt^cx^ay^b:a,b,c,d\geq0\}
\]
is a $\kk$-basis of $A$, and $A$ is a domain.
\end{lemma}

\begin{proof}
We first construct the asserted extension rather than infer it from a generic
parameter.  On $\kk[z,t]$, let
\[
 \sigma_1(z)=z,\qquad \sigma_1(t)=qt.
\]
On $S=\kk[z,t][x;\sigma_1]$, define
\[
 \begin{array}{c|ccc}
       &z&t&x\\ \hline
 \sigma_2&z&q^2t&qx\\
 \delta_2&0&z&-qt.
 \end{array}
\]
The map $\sigma_2$ preserves $xt=qtx$.  For the associated
$\sigma_2$-derivation, the only nontrivial compatibility is the same
relation, and
\[
 \delta_2(xt)=qxz-qt^2,
 \qquad
 \delta_2(qtx)=qzx-qt^2.
\]
These are equal because $z$ is central in $S$; the relations involving $z$
are immediate.  Thus the iterated Ore extension exists and has precisely
the rules \eqref{eq:ore-rules}.  Since the two $\sigma_i$ are automorphisms
and the base ring is a domain, the displayed PBW monomials form a basis and
the extension is a domain.

It remains to identify this extension with the algebra defined by
\eqref{eq:specialized-serre}.  Inside the Ore extension, the last relation in
\eqref{eq:ore-rules} gives
\[
 t=xy-q^2yx,
\]
and the preceding one gives $z=yt-q^2ty$.  The rule $xt=qtx$ is equivalent
to the first relation in \eqref{eq:specialized-serre}.  Moreover, two direct
reductions give
\[
 y^2x=q^2xy^2+qty-qz,
 \qquad
 y^3x=xy^3.
\]
Hence $x,y$ satisfy both specialized Serre relations, and they generate the
Ore extension because they recover $t$ and $z$.

Conversely, define $t=xy-q^2yx$ and $z=yt-q^2ty$ inside the algebra given by
\eqref{eq:specialized-serre}.  Its first relation gives $xt=qtx$, while the
definitions give the last two rules in \eqref{eq:ore-rules}.  Direct
expansion using $xt=qtx$ yields
\begin{equation}
 yz-zy=q^2(xy^3-y^3x),
 \label{eq:yz}
\end{equation}
so the second relation in \eqref{eq:specialized-serre} gives $yz=zy$.
Also,
\[
 xz=t^2+q^2yxt-txy=zx.
\]
Thus $z$ commutes with $x$ and $y$, and therefore with $t$.  All the Ore
rules hold.  The homomorphisms just constructed fix $x,y$ and recover $t,z$,
so they are inverse.
\end{proof}

\section{The full center}

We retain the notation $x,y,t,z$ of Lemma~\ref{lem:ore} and put
\begin{equation}
 s=xyt+\frac{xz}{q^2-1}+\frac{t^2}{q^4-1}.
 \label{eq:s}
\end{equation}
Thus $s=z_1$ under \eqref{eq:e3}--\eqref{eq:z1}.

\begin{lemma}\label{lem:central-generators}
The elements $x^3,y^3,t^3,z,s$ are central in $A$.
\end{lemma}

\begin{proof}
Multiplying the first equation of \eqref{eq:specialized-serre} on the left
and right by $x$ and subtracting proves $[x^3,y]=0$.  The second equation is
$[y^3,x]=0$, so $x^3$ and $y^3$ are central.

From \eqref{eq:ore-rules}, induction gives
\[
 yt^n=q^{2n}t^ny+(1+q^2+\cdots+q^{2(n-1)})zt^{n-1}.
\]
At $n=3$ the geometric sum is zero and $q^6=1$.  Together with
$xt=qtx$, this proves that $t^3$ is central.  Centrality of $z$ was proved
in Lemma~\ref{lem:ore}.

For $s$, set
\[
 \alpha=\frac1{q^2-1},\qquad
 \beta=\frac1{q-1},\qquad
 \gamma=1+\alpha=-\beta.
\]
Normal ordering \eqref{eq:s} gives
\[
 s=txy+\gamma zx+\beta t^2.
\]
The identity $1+q+q^2=0$ implies
\[
 1+q\gamma=\gamma,
 \qquad -1+q\beta=\beta,
 \qquad \beta-q^2=\beta q^2,
 \qquad \gamma+\beta=0.
\]
Reduction with \eqref{eq:ore-rules} now gives
\begin{align*}
 st-ts&=(1+q\gamma-\gamma)ztx=0,\\
 sx-xs&=(\beta-q^2-\beta q^2)t^2x=0,\\
 ys-sy&=(-1+q\beta-\beta)t^2y
       +(1+q\gamma-\gamma)zxy-q(\gamma+\beta)zt=0.
\end{align*}
The element $s$ also commutes with the central element $z$.  Since
$x,y$ generate $A$, $s$ is central.
\end{proof}

\begin{proposition}\label{prop:center}
The full center is
\[
 Z(A)=\kk[x^3,y^3,t^3,z,s].
\]
\end{proposition}

\begin{proof}
Filter $A$ by total exponent in the PBW generators $z,t,x,y$ from
Lemma~\ref{lem:ore}.  The terms $z$ in $yt=q^2ty+z$ and $t$ in
$yx=qxy-qt$ have lower filtered degree.  Therefore $\gr A$ is the quantum
affine space with $Z$ central and
\begin{equation}
 XT=qTX,
 \qquad YT=q^2TY,
 \qquad YX=qXY,
 \label{eq:gr-relations}
\end{equation}
where capital letters denote principal symbols.

Consider the PBW monomial $T^cX^aY^bZ^d$.  Commuting it in turn with
$T,X,Y$ shows that it is central exactly when
\begin{equation}
 a+2b\equiv0,\qquad b-c\equiv0,\qquad a-c\equiv0
 \pmod 3.
 \label{eq:center-congruences}
\end{equation}
These conditions are equivalent to $a\equiv b\equiv c\pmod3$; the exponent
$d$ is unrestricted.  PBW independence shows as well that a central
polynomial is a linear combination of precisely these central monomials.

The nonnegative exponent semigroup described by
\eqref{eq:center-congruences} is generated by the three coordinate cubes,
the free $Z$ coordinate, and the common-residue vector $(1,1,1,0)$ (with
coordinates ordered as the exponents of $X,Y,T,Z$).  Indeed, subtract the
common residue $r\in\{0,1,2\}$ and divide the first three coordinates by
three; the case $r=2$ is the square of the common-residue generator, up to a
nonzero quantum scalar.  Hence
\[
 Z(\gr A)=\kk[X^3,Y^3,T^3,Z,TXY].
\]

By Lemma~\ref{lem:central-generators}, the first four generators have
central lifts $x^3,y^3,t^3,z$.  The highest filtered symbol of $s$ is the
nonzero monomial $TXY$.  If $C\in Z(A)$, then its highest filtered symbol is
central in $\gr A$.  Choose a homogeneous polynomial in
$X^3,Y^3,T^3,Z,TXY$ equal to that symbol and substitute the five central
lifts.  Subtracting the result from $C$ strictly lowers filtered degree and
preserves centrality.  Induction on filtered degree terminates at the
scalars and proves the claimed equality.
\end{proof}

\section{Automorphisms fixing the center}

We now use the original connected grading
\[
 \deg x=\deg y=1,\qquad \deg t=2,\qquad \deg z=3.
\]
All relations in \eqref{eq:ore-rules} are homogeneous for this grading.
Because $A$ is a graded domain, the highest degree of a product of nonzero
elements is the sum of their highest degrees.

\begin{proposition}\label{prop:ozone}
Every $\kk$-algebra automorphism of $A$ that fixes $Z(A)$ pointwise is the
identity.
\end{proposition}

\begin{proof}
Let $\varphi\in\Aut_{Z(A)\text{-alg}}(A)$.  By
Proposition~\ref{prop:center},
\[
 \varphi(x)^3=x^3,
 \qquad
 \varphi(y)^3=y^3.
\]
Degree additivity in the domain gives $\deg\varphi(x)=1$.  The degree-zero
part of $\varphi(x)^3=x^3$ shows that the scalar term of $\varphi(x)$ is
zero.  Thus
\[
 \varphi(x)=a x+b y
\]
for some $a,b\in\kk$.  The quotient by the two-sided ideal $(x)$ is exactly
$\kk[y]$.  Reducing the cube equality modulo $(x)$ gives $b^3y^3=0$, hence
$b=0$; then $a^3=1$.  The same argument, using
$A/(y)\cong\kk[x]$, gives scalars $\alpha,\beta\in\kk^*$ such that
\begin{equation}
 \varphi(x)=\alpha x,
 \qquad
 \varphi(y)=\beta y,
 \qquad
 \alpha^3=\beta^3=1.
 \label{eq:diagonal}
\end{equation}

The definitions of $t$ and $z$ then give
\[
 \varphi(t)=\alpha\beta t,
 \qquad
 \varphi(z)=\alpha\beta^2 z.
\]
The nonzero central PBW element $z$ is fixed, so
$\alpha\beta^2=1$.  Each term of \eqref{eq:s} scales by
$\alpha^2\beta^2$, and the nonzero element $s$ is fixed, so
$\alpha^2\beta^2=1$.  Dividing these two equations gives $\alpha=1$.
Now $\beta^2=\beta^3=1$, whence $\beta=1$.  Therefore $\varphi$ fixes the
generators $x,y$ and is the identity.
\end{proof}

Proposition~\ref{prop:center} and Proposition~\ref{prop:ozone} prove
Theorem~\ref{thm:main}.

\section{Exact checks and assistance disclosure}

An accompanying Python checker performs exact arithmetic in
$\mathbb{Q}[q]/(q^2+q+1)$ and normal-orders words using
\eqref{eq:ore-rules}.  It checks the specialized Serre identities, recovery
of $t$ and $z$, and the commutators of $x^3,y^3,t^3,z,s$ with all four PBW
generators.  It also checks the equivalence in
\eqref{eq:center-congruences} for $0\leq a,b,c<10$.  The finite check is
corroboration only: the Ore-extension construction and the infinite
semigroup/leading-term argument are proved above.  Exact residual checks
with a verification implementation separated from discovery are also used
in certified algebraic searches such as \cite{Zeng2026UMTC}; here the
program has the narrower role just stated.

Generative AI assisted with exploration of rewrite identities, construction
of the exact checker, independent reconstruction of the mathematical
argument in a separate model context, and drafting and editing the
manuscript.  Such model-context reconstruction is not peer review.  All
claims retained in the article are supported by the displayed symbolic
arguments; the author remains responsible for the result and exposition.

\bibliographystyle{abbrv}
\bibliography{references}

\end{document}
